\chapter{MPV (Mean Platelet Volume)}

\section*{What Is This Test?}
The mean platelet volume (MPV) is the average size of circulating platelets, expressed in femtoliters (fL). It is a calculated parameter derived from the platelet volume distribution histogram generated by automated hematology analyzers and is reported as part of the complete blood count. The MPV reflects platelet production activity in the bone marrow: larger platelets are generally younger and more metabolically active, having recently been released from megakaryocytes. The MPV is the platelet analog of the MCV for red blood cells and provides clinically useful information about the mechanism of thrombocytopenia and the risk of thrombotic events \cite{tierney2023current}. In thrombocytopenia, an elevated MPV suggests peripheral destruction or consumption with compensatory increased thrombopoiesis (as in immune thrombocytopenic purpura), while a low MPV suggests impaired platelet production (as in aplastic anemia or chemotherapy-induced myelosuppression) \cite{wallach2022interpretation}. The MPV is also an emerging marker of increased thrombotic and cardiovascular risk, with large, hyperfunctional platelets associated with acute coronary syndromes, stroke, and venous thromboembolism.

\section*{Alternative Names}
MPV; Mean Platelet Volume; Average Platelet Size; Platelet Volume Index; Mean Thrombocyte Volume

\section*{Principle}
The MPV is calculated from the platelet volume distribution histogram generated by automated hematology analyzers. In impedance-based systems (Beckman Coulter DxH, Sysmex XN), platelets pass through a small electrical aperture and generate resistance pulses proportional to their volume. The analyzer records the volume distribution of thousands of platelets and calculates the mean, which is reported as the MPV. Optical platelet analysis using laser light scatter provides additional volume information. The MPV is sensitive to the anticoagulant used and the time elapsed between blood collection and analysis. EDTA causes platelet swelling over time, with the MPV increasing by 10--20\% within the first 2 hours of collection and continuing to rise with prolonged storage. For this reason, MPV measurements are most accurate when performed within 2 hours of blood collection \cite{tietz2012textbook}. Sodium citrate causes less platelet swelling and provides more consistent MPV results, but is not routinely used for CBC analysis \cite{bain2017dacie}. The immature platelet fraction (IPF), measured by fluorescent flow cytometry, correlates with the MPV and provides additional information about thrombopoietic activity, but is a separate parameter.

\section*{History}
The MPV became available as a clinical parameter in the 1970s and 1980s with the development of automated hematology analyzers that generated platelet volume distribution histograms. Early studies established the relationship between platelet size and thrombopoietic activity, demonstrating that large platelets (increased MPV) were associated with accelerated thrombopoiesis in conditions of increased platelet destruction. The recognition that MPV is a marker of platelet activation and a predictor of thrombotic risk emerged in the 1990s--2000s from studies showing that patients with acute myocardial infarction and ischemic stroke had higher MPV values. The MPV is now a standard parameter reported with every CBC on modern analyzers.

\section*{Why This Test Is Ordered}
\begin{itemize}
  \item Differentiation of thrombocytopenia mechanisms --- elevated MPV with thrombocytopenia suggests peripheral destruction (ITP, DIC, TTP) with compensatory increased thrombopoiesis; low or normal MPV with thrombocytopenia suggests hypoproduction (aplastic anemia, chemotherapy, marrow infiltration)
  \item Diagnosis of immune thrombocytopenic purpura (ITP) --- elevated MPV with thrombocytopenia, giant platelets on peripheral smear, and otherwise normal CBC
  \item Identification of congenital platelet disorders --- Bernard-Soulier syndrome (markedly elevated MPV with giant platelets), Wiskott-Aldrich syndrome (low MPV with microthrombocytes)
  \item Prediction of thrombotic risk --- elevated MPV is independently associated with increased risk of myocardial infarction, ischemic stroke, and venous thromboembolism \cite{sansanayudh2014mean}
  \item Monitoring bone marrow recovery after chemotherapy or stem cell transplantation --- rising MPV and IPF precede platelet count recovery
  \item Assessment of myeloproliferative neoplasms --- elevated MPV with thrombocytosis (essential thrombocythemia, polycythemia vera)
\end{itemize}

\section*{Primary vs Secondary Test}
The MPV is a secondary calculated parameter reported with the CBC and serves as a supportive indicator of platelet production activity and thrombotic risk.

\section*{How This Test Is Performed}
The MPV is derived automatically from a venous blood sample collected in a lavender-top (EDTA) tube as part of a standard complete blood count. Approximately 2--5 mL of blood is drawn from a peripheral vein, and the tube is gently inverted 8--10 times. The sample is processed on an automated hematology analyzer, which generates a platelet volume histogram and calculates the mean \cite{fischbach2023manual}. No separate sample or additional testing is required beyond the CBC. The venipuncture takes less than 5 minutes, and automated analysis is completed within 1--2 minutes.

\section*{What Preparation Is Needed}
No special preparation is required. Fasting is not necessary. However, patients should avoid vigorous exercise immediately before blood collection, as exercise can increase platelet count and alter MPV. Medications that may affect MPV, particularly antiplatelet agents (aspirin, clopidogrel), anticoagulants, erythropoietin, thrombopoietin receptor agonists (eltrombopag, romiplostim), and chemotherapy, should be disclosed. Smoking and diabetes mellitus are associated with elevated MPV.

\section*{Contraindications and Precautions}
\begin{itemize}
  \item No absolute contraindications. Critical pre-analytical factors affecting MPV accuracy include: EDTA anticoagulant causes progressive platelet swelling --- MPV increases by 10--20\% within 2 hours and continues to rise over 24 hours, so samples should be analyzed within 1--2 hours of collection for reliable MPV; sodium citrate provides more stable MPV measurements but is not routinely used for CBC; prolonged storage at room temperature (>4 hours) causes platelet activation, swelling, and disintegration with unreliable MPV; refrigerated storage (2--8\textdegree{}C) causes cold-induced platelet shape change and activation; inadequate filling of the collection tube (excess EDTA relative to blood) causes platelet swelling and falsely elevated MPV; hemolyzed, lipemic, or clotted specimens should be rejected; very high platelet counts (>1,000,000/$\mu$L) may exceed the linear range of volume measurement; giant platelets (as in Bernard-Soulier syndrome, ITP) may exceed the upper gate threshold and be excluded, underestimating the true MPV; the MPV normal range is analyzer-specific and may differ significantly between different manufacturer platforms (Sysmex, Beckman Coulter, Siemens) --- each laboratory should establish its own reference range for the specific instrument used.
\end{itemize}
\section*{Risks}
Risks are those of routine venipuncture: mild pain or stinging at the needle site, bruising, hematoma, lightheadedness, and rarely, infection or nerve injury.

\section*{What Happens After the Test}
MPV results are available within 30--60 minutes as part of the CBC. The MPV is interpreted in the context of the platelet count, platelet distribution width (PDW), and immature platelet fraction (IPF), when available. In thrombocytopenia with elevated MPV, the smear is examined for giant platelets and confirmation of increased thrombopoiesis; ITP, DIC, or TTP workup may follow. In thrombocytopenia with low MPV, bone marrow evaluation for impaired production is considered. In patients with cardiovascular risk factors, an elevated MPV may prompt more aggressive risk factor modification, though MPV is not yet widely incorporated into formal risk stratification algorithms.

\section*{Understanding Results}

\subsection*{Below Normal (Low MPV: <7.0 fL)}
\begin{itemize}
  \item \textbf{Bone marrow failure states:} Aplastic anemia (pancytopenia with low MPV and no compensatory thrombopoiesis), chemotherapy-induced myelosuppression, radiation therapy, myelodysplastic syndromes.
  \item \textbf{Wiskott-Aldrich syndrome:} X-linked recessive disorder characterized by microthrombocytopenia (low MPV with small platelets), eczema, and recurrent infections.
  \item \textbf{Hypersplenism:} Sequestration of larger, younger platelets in the spleen, leaving smaller older platelets in circulation.
  \item \textbf{Reactive thrombocytosis:} Platelets produced in response to infection, inflammation, or iron deficiency may be smaller (low MPV), reflecting non-clonal thrombopoiesis with normal-size platelets.
\end{itemize}

\subsection*{Above Normal (High MPV: >10.5 fL)}
\begin{itemize}
  \item \textbf{Immune thrombocytopenic purpura (ITP):} Elevated MPV with thrombocytopenia and giant platelets on smear, reflecting accelerated thrombopoiesis to compensate for immune-mediated peripheral destruction \cite{kaushansky2021williams}.
  \item \textbf{Disseminated intravascular coagulation (DIC) and thrombotic thrombocytopenic purpura (TTP):} Increased platelet consumption with compensatory release of large, young platelets.
  \item \textbf{Bernard-Soulier syndrome:} Autosomal recessive disorder with markedly elevated MPV, giant platelets (approaching or exceeding red blood cell size), and bleeding due to deficiency of the platelet glycoprotein Ib-IX-V complex (von Willebrand factor receptor).
  \item \textbf{Myeloproliferative neoplasms:} Essential thrombocythemia and polycythemia vera often have elevated MPV with abnormal megakaryocyte proliferation.
  \item \textbf{May-Hegglin anomaly:} Autosomal dominant disorder with giant platelets (elevated MPV), thrombocytopenia, and D\"{o}hle bodies in neutrophils.
  \item \textbf{Acute coronary syndrome and myocardial infarction:} Elevated MPV has been shown in multiple studies to be independently associated with increased risk of myocardial infarction and adverse outcomes post-MI, likely reflecting larger, more reactive platelets that aggregate more readily \cite{chu2010mean}.
  \item \textbf{Ischemic stroke and transient ischemic attack:} Elevated MPV is associated with increased risk of ischemic stroke, severity of stroke, and poor functional outcome.
  \item \textbf{Diabetes mellitus:} Hyperglycemia and insulin resistance are associated with increased MPV and heightened platelet reactivity.
  \item \textbf{Smoking:} Nicotine and tobacco-related oxidative stress increase platelet activation and MPV.
  \item \textbf{Sepsis:} Pro-inflammatory cytokines stimulate megakaryopoiesis and release of large, activated platelets.
  \item \textbf{Post-splenectomy:} Loss of splenic sequestration removes the normal quality-control function of the spleen (which removes older, smaller platelets), shifting the platelet population toward larger, younger forms.
\end{itemize}

\section*{Normal Results}
\begin{itemize}
  \item \textbf{Adults:} 7.0--10.5 fL (range varies by analyzer platform --- Beckman Coulter: 7.0--10.5 fL; Sysmex: 9.0--12.0 fL; Siemens ADVIA: 6.5--10.0 fL)
  \item \textbf{Newborns:} 7.0--10.0 fL (similar to adult range)
  \item \textbf{Children:} 7.0--10.5 fL (similar to adult range)
  \item \textbf{Note:} MPV is inversely related to platelet count in normal individuals (larger platelet populations require fewer cells to maintain platelet mass), and this physiological inverse relationship must be considered when interpreting results.
\end{itemize}

\section*{Follow-up Tests}
\begin{itemize}
  \item \textbf{Peripheral blood smear:} Essential for visual assessment of platelet size, morphology, giant platelets, platelet clumps, and associated red/white cell abnormalities
  \item \textbf{Immature platelet fraction (IPF):} Provides a direct measure of young, RNA-rich platelets newly released from the bone marrow; a more direct indicator of thrombopoietic activity than MPV alone
  \item \textbf{Bone marrow aspiration and biopsy:} When aplastic anemia, MDS, leukemia, or marrow infiltration is suspected with low MPV
  \item \textbf{ADAMTS13 activity:} When TTP is suspected (thrombocytopenia with elevated MPV and microangiopathic hemolytic anemia)
  \item \textbf{Coagulation panel (PT, PTT, fibrinogen, D-dimer):} When DIC is suspected
  \item \textbf{JAK2 V617F, CALR, MPL mutation testing:} When myeloproliferative neoplasm with thrombocytosis and elevated MPV is suspected
  \item \textbf{HIV, hepatitis C, EBV serologies:} When infection-associated thrombocytopenia is considered
  \item \textbf{Platelet aggregation studies (light transmission aggregometry):} To assess platelet function when bleeding or thrombotic disorder is suspected
  \item \textbf{Flow cytometry for platelet glycoproteins:} When Bernard-Soulier syndrome (GP Ib-IX-V deficiency) or Glanzmann thrombasthenia (GP IIb/IIIa deficiency) is suspected
\end{itemize}

\section*{References}
\bibliographystyle{unsrtnat}
\bibliography{references}

\clearpage
\section*{Regulatory Status and Authorization Date}
This chapter describes a test category, not a specific commercial product. FDA, CDSCO, EU IVDR, and MHRA approval or registration dates therefore cannot be assigned without the assay manufacturer, product name, instrument, and intended use. For laboratory-developed tests, an individual agency approval date may not exist. See the Preface for the interpretation of regulatory status.

\clearpage
